% Options for packages loaded elsewhere
\PassOptionsToPackage{unicode}{hyperref}
\PassOptionsToPackage{hyphens}{url}
\documentclass[
]{article}
\usepackage{xcolor}
\usepackage{amsmath,amssymb}
\setcounter{secnumdepth}{-\maxdimen} % remove section numbering
\usepackage{iftex}
\ifPDFTeX
  \usepackage[T1]{fontenc}
  \usepackage[utf8]{inputenc}
  \usepackage{textcomp} % provide euro and other symbols
\else % if luatex or xetex
  \usepackage{unicode-math} % this also loads fontspec
  \defaultfontfeatures{Scale=MatchLowercase}
  \defaultfontfeatures[\rmfamily]{Ligatures=TeX,Scale=1}
\fi
\usepackage{lmodern}
\ifPDFTeX\else
  % xetex/luatex font selection
\fi
% Use upquote if available, for straight quotes in verbatim environments
\IfFileExists{upquote.sty}{\usepackage{upquote}}{}
\IfFileExists{microtype.sty}{% use microtype if available
  \usepackage[]{microtype}
  \UseMicrotypeSet[protrusion]{basicmath} % disable protrusion for tt fonts
}{}
\makeatletter
\@ifundefined{KOMAClassName}{% if non-KOMA class
  \IfFileExists{parskip.sty}{%
    \usepackage{parskip}
  }{% else
    \setlength{\parindent}{0pt}
    \setlength{\parskip}{6pt plus 2pt minus 1pt}}
}{% if KOMA class
  \KOMAoptions{parskip=half}}
\makeatother
\setlength{\emergencystretch}{3em} % prevent overfull lines
\providecommand{\tightlist}{%
  \setlength{\itemsep}{0pt}\setlength{\parskip}{0pt}}
\usepackage{bookmark}
\IfFileExists{xurl.sty}{\usepackage{xurl}}{} % add URL line breaks if available
\urlstyle{same}
\hypersetup{
  pdftitle={Effective Real Time Embedded C and C++},
  hidelinks,
  pdfcreator={LaTeX via pandoc}}

\title{Effective Real Time Embedded C and \texttt{C++}}
\author{}
\date{}

\begin{document}
\maketitle

\subsection{Description}\label{description}

Many software engineers work with C and \texttt{C++} programming
languages for real time and embedded development. This course aims to
give them a fuller appreciation of the more advanced aspects of the
languages in the real time and embedded environments. Programming real
time embedded systems \texttt{make} many special demands on the software
engineer. These include understanding compile, link and run-time issues,
which are included in the course. It also covers real time programming
issues that are specific to the C and \texttt{C++} languages. Most of
the course is relevant to both C and \texttt{C++} software engineers,
though some special points relevant to \texttt{C++} and object oriented
are also presented. The course is specifically aimed at issues relevant
to real-time embedded software engineers.

\subsection{Duration}\label{duration}

4 days / 32 hours

\subsection{Intended Audience}\label{intended-audience}

The course is designed for software engineers who want to improve their
use of the C and / or \texttt{C++} programming languages in a real time,
embedded environment.

\subsection{Prerequisites}\label{prerequisites}

\begin{itemize}
\tightlist
\item
  A good grasp of the fundamentals of C and/or \texttt{C++}
\item
  A good grasp of real time and embedded is an advantage.
\end{itemize}

\subsection{Objectives}\label{objectives}

\begin{itemize}
\tightlist
\item
  understand the core concepts and principles of Effective Real Time
  Embedded C and \texttt{C++}
\item
  gain practical knowledge of Introduction
\item
  gain practical knowledge of Embedded C
\item
  gain practical knowledge of Elements of C/`\texttt{C++}
\end{itemize}

\subsection{Outline}\label{outline}

\begin{itemize}
\tightlist
\item
  Introduction
\item
  Embedded C

  \begin{itemize}
  \tightlist
  \item
    Real time C
  \item
    Embedded C
  \item
    Writing code in Kernel Space
  \end{itemize}
\item
  Elements of C/\texttt{C++}

  \begin{itemize}
  \tightlist
  \item
    Type Specifiers and Qualifiers
  \item
    Type coercion and conversion
  \item
    Struct layout
  \item
    Uses of union
  \item
    Bitfields
  \item
    Enum v \#define
  \item
    Typedef
  \item
    Unscrambling declarations (cdecl)
  \item
    Floating-point and fixed-point number systems
  \item
    IEEE Standard 754
  \item
    Importance of good structure
  \item
    Quality and style
  \end{itemize}
\item
  Effective Pointers

  \begin{itemize}
  \tightlist
  \item
    Arrays and pointers; compatibility and incompatibility
  \item
    The principle of dynamic memory allocation
  \item
    Void pointers
  \item
    Implementing basic data structures: pointers, arrays, queues,
    stacks, lists
  \item
    Function pointers: basics, callbacks, state machine
  \item
    Optimizations with pointers
  \item
    Variable sized structures
  \end{itemize}
\item
  Memory Management Architecture

  \begin{itemize}
  \tightlist
  \item
    Caching
  \item
    Hardware memory management
  \item
    Sections of a program
  \item
    Mapping of program in memory
  \item
    Dynamic memory allocation
  \end{itemize}
\item
  Intertask Communication

  \begin{itemize}
  \tightlist
  \item
    Intertask Synchronization: Race Conditions and Critical Sections
  \item
    Shared Memory
  \item
    Locks
  \item
    Mutexes
  \item
    Semaphores
  \item
    Deadlocks
  \item
    Signals
  \item
    Message Queues
  \end{itemize}
\item
  The Toolchain

  \begin{itemize}
  \tightlist
  \item
    Pre-processor compiler directives and cross compilers
  \item
    Macros v functions
  \item
    The Linker and memory sections
  \item
    Using the toolchain to your advantage
  \item
    What happens before main
  \item
    What is happening at runtime
  \item
    Estimation of stack requirements
  \end{itemize}
\item
  Optimizations

  \begin{itemize}
  \tightlist
  \item
    Tips for optimizing embedded C/\texttt{C++} code
  \end{itemize}
\item
  Writing Safer C

  \begin{itemize}
  \tightlist
  \item
    \texttt{MISRA}-C and \texttt{CWE}
  \item
    Compiler warnings
  \item
    Hard learned pitfalls
  \item
    Coding style and banned \texttt{API}
  \item
    Preparing for debug
  \end{itemize}
\item
  Hardware Programming Model

  \begin{itemize}
  \tightlist
  \item
    Interrupts - handlers and dispatchers
  \item
    Hardware and Software interrupts
  \item
    Contexts
  \end{itemize}
\item
  Timing

  \begin{itemize}
  \tightlist
  \item
    The challenge in measuring time
  \item
    Hardware and Software timers
  \item
    Watchdog timer
  \item
    Real time debug vs.~time
  \end{itemize}
\item
  Object oriented C

  \begin{itemize}
  \tightlist
  \item
    Implementing \texttt{OOP} with C (OOC)
  \item
    Mixing C and \texttt{C++}
  \end{itemize}
\item
  Inheritance and \texttt{OO} Design

  \begin{itemize}
  \tightlist
  \item
    Inheritance models
  \item
    The use of \texttt{virtual-functions}
  \end{itemize}
\item
  Embedded \texttt{C++}

  \begin{itemize}
  \tightlist
  \item
    C vs.~\texttt{C++}
  \item
    Benefits of embedded \texttt{C++} (EC++)
  \item
    Restrictions for EC++
  \end{itemize}
\item
  \texttt{C++} Resource Management

  \begin{itemize}
  \tightlist
  \item
    Managing contained members
  \item
    The Rule of the Big Three
  \item
    Garbage collection
  \item
    Smart pointers
  \item
    Resource managing classes
  \end{itemize}
\item
  Templates and Generic Programming

  \begin{itemize}
  \tightlist
  \item
    Understand the two meanings of typename
  \item
    Know how to access names in templatized base classes
  \item
    Parameter-independent code
  \item
    Traits
  \item
    Template metaprogramming
  \end{itemize}
\item
  The Standard Template Library

  \begin{itemize}
  \tightlist
  \item
    \texttt{STL} Components
  \item
    Containers
  \item
    Iterators
  \item
    Algorithms
  \item
    Using emplace
  \item
    Manipulating Algorithms
  \item
    Container Elements
  \item
    Extending the \texttt{STL}
  \end{itemize}
\item
  Input/Output Using Stream Classes

  \begin{itemize}
  \tightlist
  \item
    Common Background of \texttt{I/O} Streams
  \item
    Fundamental Stream Classes and Objects
  \item
    Standard Stream Operators \textless\textless{} and
    \textgreater\textgreater{}
  \item
    Standard Input/Output Functions
  \item
    Manipulators
  \item
    Formatting
  \item
    Internationalization
  \item
    File Access
  \item
    Connecting Input and Output Streams
  \item
    Stream Classes for Strings
  \item
    Input/Output Operators for User-Defined Types
  \item
    The Stream Buffer Classes
  \item
    Performance Issues
  \end{itemize}
\item
  Peripherals (extra)

  \begin{itemize}
  \tightlist
  \item
    General Purpose Input/Output (\texttt{GPIO})
  \item
    Universal Asynchronous Receiver-Transmitter (\texttt{UART})
  \item
    Analog-to-Digital Converter (ADC)
  \item
    Serial Peripheral Interface (\texttt{SPI})
  \item
    Inter-Integrated Circuit (\texttt{I2C})
  \item
    Direct Memory Access (\texttt{DMA})
  \item
    Nested Vector Interrupt Controller (NVIC)
  \item
    General Purpose Timers (TIM)
  \item
    System Tick Timer (SysTick)
  \end{itemize}
\end{itemize}

\subsection{Copyright}\label{copyright}

Mark Veltzer
\href{mailto:mark.veltzer@gmail.com}{\nolinkurl{mark.veltzer@gmail.com}},
© 2026

\end{document}
